\section{Changing the observation declaration}
\label{sim:joint-domain-changes}

The fixed-domain result raises a practical question when a modeller revises the observation declaration. A message may acquire a new label, or a formerly pooled observation may be separated into more informative messages. Both changes alter the encoded decision keys, but they have different consequences. A new spelling can preserve the same choices exactly. A new distinction can let a manager choose different actions where the old declaration required the same action. The comparison must say which change has occurred before reusing a computed policy set.

\subsection{A different label must preserve the entire decision key}

Consider replacing every first-message label with a distinct new label and doing the same for the second messages. The substitution must be global across source heads: two equal old messages must remain equal, and distinct old messages must remain distinct. A substitution chosen separately for each source could reveal the source's identity and therefore change the information available to the manager.

The second-stage key contains the first message and the manager's first action as well as the current message. Consequently, renaming the first message also changes its occurrence inside second-stage memory. The joint comparison generates this full key correspondence from all original financial and observation-head triples, before inspecting completed source outcomes. If the correspondence is a bijection, the old total action assignment can be permuted to the new key order. Inverting the permutation recovers every assignment.

The refusal evidence travels with the assignment. A rejected rule keeps the same financial-row identity, but the checker must confirm that the renamed rule still takes the actions that row refuses. The transported surviving population agrees with fresh evaluation under the new declaration, and the source supports, pending clock regions and unresolved source outcomes remain unchanged. An overflowing policy domain stays unassessed; label equivalence does not remove its computational limit.\footnote{The bounded joint comparison in \path{data/accounting-master/joint-domain-transport-20261027/} checks 34 comparisons across nine declarations and 16 alteration controls. Eight comparisons retain unassessed domain-budget results. The assessed comparisons check full populations, preserved refusal witnesses and inverse roundtrips, including the first message inside second-stage memory.}

\subsection{A new distinction permits an old rule to ignore it}

Now suppose a refined declaration distinguishes messages that a coarse declaration pools. The financial rows, action meanings, source chronology and decision dates are unchanged. Let $D_c$ and $D_f$ be the coarse and refined sets of full decision keys. A surjective map $e:D_f\to D_c$ erases the new distinctions. It must also erase the first-message distinction wherever that message occurs in second-stage memory. Surjectivity ensures that every coarse decision key still has a refined counterpart.

A coarse rule $\pi$ assigns an action to each key in $D_c$. Its lift $\pi^{\uparrow}$ on $D_f$ ignores the added distinction: at a refined key $z$, it takes the coarse action at $e(z)$. Let $C_c$ and $C_f$ be the candidate populations for the same selected source pairs under the two declarations. With unchanged financial refusal conditions, the checked relation is
\begin{equation}
 \pi^{\uparrow}(z)=\pi(e(z)),\qquad
 \{\pi^{\uparrow}:\pi\in C_c\}
 =\{\rho\in C_f:e(z)=e(z')\Rightarrow\rho(z)=\rho(z')\}.
 \label{eq:joint-observation-lift}
\end{equation}
The right-hand side selects refined rules that take the same action whenever the new messages collapse to the same coarse key. To see why the equality holds, substitute the coarse action at $e(z)$ for each refined action in a refused financial row. Erasing that row's decision keys produces exactly the coarse refusal condition. A lifted coarse candidate therefore passes every refined row. Conversely, a refined candidate constant on each erasure class defines a unique coarse rule, and the same substitution shows that it passes the coarse rows. The equality concerns the same finite decision problem under two declared information restrictions.

In the between-date correction case, restoring the original first-message distinctions from an artificially collapsed first message changes the candidate population from zero to 16 total rules, with four additional reached-row signatures. Restoring the second-message distinctions changes two total candidates to 16, with three additional signatures. A signature records the financial rows reached by a rule's actions; it is not a probability distribution or an observed company path. The comparison verifies the reached rows of each individual lifted rule, so equality of aggregate signature sets cannot conceal a wrong pairing.

The distinction matters particularly when required sources are unresolved. In the late-conflict example, refinement expands eight total assignments to 2,048, yet no financial row is selected and no additional reached-row signature appears. Those assignments do not establish feasible company behavior. Inconsistent clocks produce no assessed policy population at all. If a refined domain exceeds the key allowance, its policy comparison remains unassessed even when the coarse domain has been evaluated.\footnote{\path{data/accounting-master/joint-observation-refinement-20261027/} retains 46 comparisons across 12 declaration pairs, with 14 alteration controls. Nine comparisons remain unassessed. Assessed cases verify full lifted populations, constant-class equality, erased refusal conditions and per-rule reached rows. First-stage, second-stage and simultaneous refinements are separate cases; exact renaming is an equality control.}

\enlargethispage{2\baselineskip}
These are conditional comparisons of supplied observation declarations. The extra distinctions must still be justified as information available to the actor. Their finite effects on permitted actions do not measure utility or predictive accuracy. Both relations also retain one financial source view across the two decisions; allowing a financial correction to change that view requires an account of what was believed earlier and which economic history continues afterward.
